% GENERATED FILE. Edit canonical lessons, then run npm run generate:books.
% canonical-source-hash: fnv1a64:5988410e39ec3bd8
% canonical-lessons: GU-C120-ran, GU-C120-vijli, GU-C120-tophan, GU-C120-dharti, GU-C120-cadni

\chapter{Desert, Lightning, Storm, Earth, Moonlight}
\label{ch:gu-desert-lightning-storm-earth-moonlight}
\hlchaptermodality{120}

\begin{chapteropening}
\textbf{By the end of this chapter:} \emph{I can say a desert, lightning, a storm, the earth and moonlight.}

Complete the last lesson of chapter 120: i can say a desert, lightning, a storm, the earth and moonlight.
\end{chapteropening}

\section[raṇ]{\gu{રણ} (raṇ) — a desert}

\label{lesson:GU-C120-ran}



\begin{quote}
Before the new one: say the Gujarati for a valley, then the Gujarati for a hill.
\end{quote}

\subsection*{You'll want to know: \gu{રણ}}
\textbf{\gu{રણ}} (\emph{raṇ}) — \textquotedblleft{}a desert\textquotedblright{}.

\begin{grammarlens}[title={What you've built}]
One more word for this chapter.
\end{grammarlens}

\subsection*{Your turn}
\begin{itemize}
\raggedright
  \item \emph{Say it:} \emph{raṇ}
  \item \emph{Say it:} \emph{raṇ}, once more
  \item \emph{Say it:} say \emph{raṇ}
  \item \emph{Recall:} say the Gujarati for a valley, then the Gujarati for a hill, then say \emph{raṇ} again
\end{itemize}

\begin{tcolorbox}[breakable,colback=teal!4,colframe=teal!35!black,title={Before you move on}]
What does \gu{રણ} mean? (\textquotedblleft{}A desert\textquotedblright{}.) Say it once more.
\end{tcolorbox}
\section[vījḷī]{\gu{વીજળી} (vījḷī) — lightning, electricity}

\label{lesson:GU-C120-vijli}



\begin{quote}
Before the new one: say the Gujarati for a hill, then the Gujarati for a desert.
\end{quote}

\subsection*{You'll want to know: \gu{વીજળી}}
\textbf{\gu{વીજળી}} (\emph{vījḷī}) — \textquotedblleft{}lightning, electricity\textquotedblright{}.

\begin{grammarlens}[title={What you've built}]
One more word for this chapter.
\end{grammarlens}

\subsection*{Your turn}
\begin{itemize}
\raggedright
  \item \emph{Say it:} \emph{vījḷī}
  \item \emph{Say it:} \emph{vījḷī}, once more
  \item \emph{Say it:} say \emph{vījḷī}
  \item \emph{Recall:} say the Gujarati for a hill, then the Gujarati for a desert, then say \emph{vījḷī} again
\end{itemize}

\begin{tcolorbox}[breakable,colback=teal!4,colframe=teal!35!black,title={Before you move on}]
What does \gu{વીજળી} mean? (\textquotedblleft{}Lightning, electricity\textquotedblright{}.) Say it once more.
\end{tcolorbox}
\section[tophān]{\gu{તોફાન} (tophān) — a storm}

\label{lesson:GU-C120-tophan}



\begin{quote}
Before the new one: say the Gujarati for a desert, then the Gujarati for lightning.
\end{quote}

\subsection*{You'll want to know: \gu{તોફાન}}
\textbf{\gu{તોફાન}} (\emph{tophān}) — \textquotedblleft{}a storm\textquotedblright{}.

\begin{grammarlens}[title={What you've built}]
One more word for this chapter.
\end{grammarlens}

\subsection*{Your turn}
\begin{itemize}
\raggedright
  \item \emph{Say it:} \emph{tophān}
  \item \emph{Say it:} \emph{tophān}, once more
  \item \emph{Say it:} say \emph{tophān}
  \item \emph{Recall:} say the Gujarati for a desert, then the Gujarati for lightning, then say \emph{tophān} again
\end{itemize}

\begin{tcolorbox}[breakable,colback=teal!4,colframe=teal!35!black,title={Before you move on}]
What does \gu{તોફાન} mean? (\textquotedblleft{}A storm\textquotedblright{}.) Say it once more.
\end{tcolorbox}
\section[dhartī]{\gu{ધરતી} (dhartī) — the earth, the ground}

\label{lesson:GU-C120-dharti}



\begin{quote}
Before the new one: say the Gujarati for lightning, then the Gujarati for a storm.
\end{quote}

\subsection*{You'll want to know: \gu{ધરતી}}
\textbf{\gu{ધરતી}} (\emph{dhartī}) — \textquotedblleft{}the earth, the ground\textquotedblright{}.

\begin{grammarlens}[title={What you've built}]
One more word for this chapter.
\end{grammarlens}

\subsection*{Your turn}
\begin{itemize}
\raggedright
  \item \emph{Say it:} \emph{dhartī}
  \item \emph{Say it:} \emph{dhartī}, once more
  \item \emph{Say it:} say \emph{dhartī}
  \item \emph{Recall:} say the Gujarati for lightning, then the Gujarati for a storm, then say \emph{dhartī} again
\end{itemize}

\begin{tcolorbox}[breakable,colback=teal!4,colframe=teal!35!black,title={Before you move on}]
What does \gu{ધરતી} mean? (\textquotedblleft{}The earth, the ground\textquotedblright{}.) Say it once more.
\end{tcolorbox}
\section[cā̃dnī]{\gu{ચાંદની} (cā̃dnī) — moonlight}

\label{lesson:GU-C120-cadni}



\begin{quote}
Before the new one: say the Gujarati for a storm, then the Gujarati for the earth.

In \emph{je … te …}, which half describes, and which names the one meant? (\emph{Je} describes; \emph{te} names.)
\end{quote}

\subsection*{You'll want to know: \gu{ચાંદની}}
\textbf{\gu{ચાંદની}} (\emph{cā̃dnī}) — \textquotedblleft{}moonlight\textquotedblright{}.

\begin{grammarlens}[title={What you've built}]
One more word for this chapter.
\end{grammarlens}

\subsection*{Your turn}
\begin{itemize}
\raggedright
  \item \emph{Say it:} \emph{cā̃dnī}
  \item \emph{Say it:} \emph{cā̃dnī}, once more
  \item \emph{Say it:} say \emph{cā̃dnī}
  \item \emph{Recall:} say the Gujarati for a storm, then the Gujarati for the earth, then say \emph{cā̃dnī} again
\end{itemize}

\begin{tcolorbox}[breakable,colback=teal!4,colframe=teal!35!black,title={Before you move on}]
What does \gu{ચાંદની} mean? (\textquotedblleft{}Moonlight\textquotedblright{}.) Say it once more.
\end{tcolorbox}
